\documentclass[11pt,a4paper]{article}

% --- Language & Font Setup ---
\usepackage[utf8]{inputenc}
\usepackage[T1]{fontenc}
\usepackage[ngerman]{babel}
\usepackage[default]{lato}                    % modern, ATS-safe sans-serif
\usepackage{iftex}
\setlength{\emergencystretch}{3em}
\hyphenpenalty=1000
\ifxetex
  \usepackage{microtype}
  \providecommand{\textls}[2][150]{{\addfontfeature{LetterSpace=15}#2}}
  \renewcommand{\textls}[2][150]{{\addfontfeature{LetterSpace=15}#2}}
\else
  \usepackage[letterspace=150]{microtype}
\fi
\usepackage{xcolor}

% --- Design Colors ---
\definecolor{accent}{HTML}{1F3864}      % dark navy (name, section rules, company)
\definecolor{accentlight}{HTML}{4A6FA5} % lighter blue (role, bold labels)
\definecolor{datecolor}{HTML}{595959}   % neutral gray (dates)

% --- Layout & Margins ---
\usepackage[left=2.1cm, right=2.1cm, top=1.7cm, bottom=1.7cm]{geometry}
\usepackage{titlesec}
\usepackage{enumitem}
\usepackage{hyperref}
\usepackage{color}
\usepackage{fontspec}

\linespread{1.05}

% --- ATS-Friendly Formatting ---
\pagestyle{empty}
\urlstyle{same}
\hypersetup{
    colorlinks=true,
    linkcolor=accent,
    urlcolor=accent,
    pdfborder={0 0 0},
    pdfauthor={Khaled Swaidan},
    pdftitle={Lebenslauf -- Khaled Swaidan},
    pdfsubject={Mobile Engineer -- Flutter, iOS/Android}
}

% --- Section headings: bold small caps, navy, thin rule underneath ---
\titleformat{\section}
{\normalfont\large\bfseries\scshape\color{accent}}
{}{0pt}{}
[{\vspace{2pt}{\color{accent}\titlerule[0.6pt]}}]

\titlespacing*{\section}{0pt}{14pt}{6pt}

% --- Bullet lists: colored bullet, more air ---
\setlist[itemize]{
    leftmargin=1.2em,
    labelsep=0.7em,
    label=\textcolor{accent}{\textbullet},
    itemsep=3pt,
    topsep=3pt,
    parsep=0pt,
    partopsep=0pt
}

% --- Paragraph spacing ---
\setlength{\parindent}{0pt}
\setlength{\parskip}{2pt}
\raggedbottom

% --- Custom commands ---

% Header: name (letter-spaced navy) / role (light blue) / contact (gray) + rule
\newcommand{\cvheader}[3]{%
    \begin{center}
    {\fontsize{24}{28}\selectfont\bfseries\color{accent}\textls[150]{#1}}\\[5pt]
    {\normalsize\color{accentlight}#2}\\[7pt]
    {\small\color{datecolor}#3}
    \end{center}
    \vspace{3pt}
    {\color{accent}\rule{\linewidth}{0.9pt}}
}

% Entry: role (bold black) + date (gray, right) / company (italic navy)
\newcommand{\entry}[3]{%
    \vspace{6pt}%
    \noindent\textbf{#1} \hfill {\small\color{datecolor}#2}\\[2pt]
    {\itshape\color{accent}#3}\par
}

\newcommand{\entrysimple}[2]{%
    \vspace{6pt}%
    \noindent\textbf{#1} \hfill {\small\color{datecolor}#2}\par
}

\begin{document}

    \cvheader
    {Khaled Swaidan}
    {AI Native Software Engineer -- Python, RAG \& LLMs | M.Sc. Informatik}
    {Schwachhausen, 28211 Bremen \textbullet\ \href{mailto:khaled.swaidan.k@gmail.com}{khaled.swaidan.k@gmail.com} \textbullet\ Verfügbar ab sofort}

    \section{Profil}
    AI Native Software Engineer (M.Sc. Informatik, Thesis 1,0) mit Fokus auf die Entwicklung produktionsreifer GenAI-Systeme. Ich integriere LLM-APIs und KI-Assistenten (ChatGPT, Claude, Copilot) als festen Bestandteil meines täglichen Workflows und übernehme die volle Verantwortung für die Qualität der generierten Outputs. Meine Stärke liegt in der Evaluierung von KI-Modellen (Benchmarking, Guardrails) und der End-to-End-Entwicklung von RAG-Pipelines und Full-Stack-Anwendungen.

    \section{Berufserfahrung}

    \entry{Softwareentwickler (Cloud \& Automation)}{02/2025--04/2026}{Rocket-Code}
    \begin{itemize}
        \item \textcolor{accentlight}{\textbf{AI-Assisted Engineering:}} Täglicher Einsatz von KI-Coding-Assistenten (ChatGPT, Claude, GitHub Copilot) zur Code-Generierung, Refactoring und Test-Erstellung; kritische Prüfung und Validierung der Ergebnisse.
        \item \textcolor{accentlight}{\textbf{Backend \& Cloud:}} Anbindung von Enterprise-Anwendungen an Azure Functions und REST-APIs; Sicherstellung stabiler Datenflüsse.
        \item \textcolor{accentlight}{\textbf{Agile \& CI/CD:}} Entwicklung von automatisierten Test-Pipelines (Git Hooks) und Mitwirkung in agilen Scrum-Teams.
    \end{itemize}

    \entry{Junior IT-Berater -- Requirements \& Evaluation}{11/2023--07/2024}{IFIB Consult}
    \begin{itemize}
        \item \textcolor{accentlight}{\textbf{Requirements Engineering:}} Analyse von Kundenanforderungen und Übersetzung in technische Spezifikationen und User Stories.
        \item \textcolor{accentlight}{\textbf{Qualitätssicherung:}} Externe Evaluationsstelle für Softwareprojekte; strukturierte Fehleranalyse und Dokumentation.
    \end{itemize}

    \section{Projekte \& KI-Erfahrung}

    \entry{SIRcraft AI Platform -- Masterarbeit}{10/2025--10/2026}{Universität Bremen | Note 1,0}
    \begin{itemize}
        \item \textcolor{accentlight}{\textbf{RAG \& Agentic Systems:}} Aufbau einer End-to-End RAG-Plattform mit LangChain und ChromaDB zur Anbindung interner Wissensdatenbanken; Integration von LLM-APIs (Gemini, Ollama).
        \item \textcolor{accentlight}{\textbf{Evaluation \& Guardrails:}} Entwicklung eines Python-basierten AST-Parsers zur automatisierten Validierung (Bewertung) von KI-generierten Outputs und zur Absicherung gegen Halluzinationen.
        \item \textcolor{accentlight}{\textbf{Full-Stack \& Deployment:}} Entwicklung eines FastAPI-Backends und Flutter-Frontends; Containerisierung via Docker Compose für den produktiven Betrieb.
    \end{itemize}

    \section{Bildung}

    \entry{M.Sc. Informatik}{10/2023--09/2026}{Universität Bremen | Gesamtnote 1,8; Masterarbeit 1,0}
    \entry{B.Sc. Informatik}{10/2018--09/2023}{Universität Bremen}

    \section{Fachliche Kenntnisse}
    \begin{itemize}
        \item \textcolor{accentlight}{\textbf{KI \& LLMs:}} RAG-Architekturen, LangChain, LLM-APIs (OpenAI, Gemini), Prompt Engineering, Evaluierung von KI-Outputs, Guardrails (AST).
        \item \textcolor{accentlight}{\textbf{Programmierung \& Full-Stack:}} Python, FastAPI, RESTful APIs, Flutter/Dart, SQL, HTML/CSS.
        \item \textcolor{accentlight}{\textbf{Cloud \& DevOps:}} Azure, Docker, CI/CD, Git, Testautomatisierung, Agile (Scrum).
    \end{itemize}

    \section{Sprachen}
    Deutsch (C1) \quad Englisch (C1) \quad Arabisch (Muttersprache)

\end{document}